\section{The ConvergenceSweep Algorithm}
\label{sec:algorithm}

\subsection{Notation}

Let $G = (V, E, w)$ be the weighted census-tract adjacency graph for a state
$\mathcal{S}$ with $|V| = n$ vertices and $|E| = m$ edges.
Let $k$ denote the number of congressional districts apportioned to
$\mathcal{S}$ by Huntington-Hill.
A \emph{partition} $\Pi = (D_1, \ldots, D_k)$ is a $k$-way assignment of
tracts to districts satisfying the population-balance constraint
$|\mathrm{pop}(D_i) - \bar{p}| \le \delta \bar{p}$ for each $i$, where
$\bar{p} = \mathrm{pop}(\mathcal{S})/k$ and $\delta = 0.005$ (the
$\pm 0.5\%$ federal tolerance).
The \emph{edge cut} of $\Pi$ is
\[
  \mathrm{EC}(\Pi) = \sum_{\{u,v\} \in E: \text{part}(u) \ne \text{part}(v)} w(u,v).
\]
Following B.8~\citep{b8paper}, the \emph{normalised edge cut} is
\[
  \ECnorm(\Pi) = \frac{\mathrm{EC}(\Pi)}{\sqrt{\min(i, k-i)}},
\]
where $i$ ranges over all bisection levels in the recursive tree.

The METIS bisector $\mathcal{M}(G, k, s)$ takes the graph, the district count,
and an integer seed $s$, and returns a partition $\Pi$ that is a local minimum
of $\ECnorm$ under multilevel Kernighan-Lin refinement initialised by $s$.

\subsection{Content-Derived Seed}

The starting seed is derived from the census release identifier published by
the Census Bureau under the \dia{} statute.

\begin{definition}[Content-Derived Seed]
\label{def:seed}
Let $\mathtt{census\_release\_id}$ be the ASCII string identifying the census
dataset (e.g., \texttt{"2020-PL94-171-v1.0"}), and let
$\mathtt{"DIA\_SEED\_V1"}$ be the fixed protocol constant for version 1 of the
Districting Integrity Act seed protocol.
The \emph{content-derived seed} is
\[
  s_0 = \bigl\lfloor
    \mathrm{SHA\text{-}256}(\mathtt{census\_release\_id}
      \mathbin{\|} \mathtt{"DIA\_SEED\_V1"})
  \bigr\rfloor \bmod 2^{31},
\]
where $\|$ denotes ASCII concatenation and $\lfloor \cdot \rfloor$ reads the
256-bit hash value as a big-endian unsigned integer before taking the modulus.
\end{definition}

The modulus $2^{31}$ matches the range of METIS's \texttt{int} seed parameter.
The SHA-256 pre-image is a public, verifiable string; any party can recompute
$s_0$ from first principles.

\subsection{T-Stability and the Stopping Criterion}

\begin{definition}[$T$-Stability]
\label{def:tstable}
A seed index $j$ is \emph{$T$-stable} if
\[
  \ECnorm\bigl(\mathcal{M}(G, k, s_0 + j + t)\bigr)
  \;\ge\;
  \ECnorm\bigl(\mathcal{M}(G, k, s_0 + j)\bigr)
  \quad
  \text{for all } t \in \{1, 2, \ldots, T\}.
\]
That is, none of the $T$ seeds immediately following $s_0 + j$ strictly
improves the normalised edge cut achieved at seed $s_0 + j$.
\end{definition}

\begin{definition}[Convergence Tail]
\label{def:tail}
Let $j^*$ be the index of the last seed that achieves a strict improvement
in $\ECnorm$ during a sweep.
The \emph{convergence tail} of the sweep is $\tau = j_{\mathrm{stop}} - j^*$,
where $j_{\mathrm{stop}}$ is the index at which the algorithm halts.
Under a threshold of $T$, we have $\tau = T$.
\end{definition}

\subsection{Formal Pseudocode}

Algorithm~\ref{alg:cs} gives the complete \cs{} procedure.

\begin{algorithm}[t]
\caption{\cs{($G$, $k$, $T$, $\mathtt{census\_release\_id}$)}}
\label{alg:cs}
\begin{algorithmic}[1]
\State $s_0 \gets \lfloor \mathrm{SHA\text{-}256}(\mathtt{census\_release\_id}
       \,\|\, \texttt{"DIA\_SEED\_V1"}) \rfloor \bmod 2^{31}$
\State $\Pi^* \gets \mathcal{M}(G, k, s_0)$
\State $\mathrm{EC}^* \gets \ECnorm(\Pi^*)$
\State $\mathtt{no\_improve} \gets 0$
\State $j \gets 1$
\While{$\mathtt{no\_improve} < T$}
  \State $\Pi_j \gets \mathcal{M}(G, k,\; s_0 + j)$
  \State $e_j \gets \ECnorm(\Pi_j)$
  \If{$e_j < \mathrm{EC}^*$}
    \State $\mathrm{EC}^* \gets e_j$
    \State $\Pi^* \gets \Pi_j$
    \State $\mathtt{no\_improve} \gets 0$
  \Else
    \State $\mathtt{no\_improve} \gets \mathtt{no\_improve} + 1$
  \EndIf
  \State $j \gets j + 1$
\EndWhile
\State \textbf{return} $\Pi^*$
\end{algorithmic}
\end{algorithm}

The algorithm returns the partition $\Pi^*$ achieving the minimum normalised
edge cut seen across all seeds $s_0, s_0+1, \ldots, s_0 + j - 1$.

\subsection{Correctness}

\begin{proposition}[Determinism]
\label{prop:determinism}
For fixed $G$, $k$, $T$, and $\mathtt{census\_release\_id}$, \cs{} returns
the same partition $\Pi^*$ on every execution, regardless of hardware,
operating system, or parallelisation strategy applied to independent METIS
calls.
\end{proposition}

\begin{proof}
$s_0$ is determined uniquely by Definition~\ref{def:seed}.
METIS is deterministic given $s$: the multilevel coarsening and refinement
sequence is fully specified by the seed and graph structure.
The best-partition selection in Algorithm~\ref{alg:cs} uses a strict less-than
comparison; ties (equal $\ECnorm$) do not update $\Pi^*$, so the first
minimum encountered is retained.
Therefore the entire execution path is determined by $(G, k, T, s_0)$.
\end{proof}

\begin{remark}
Proposition~\ref{prop:determinism} assumes that METIS itself is compiled with
the same floating-point model on all platforms.
The \texttt{redist} binary ships a vendored METIS with determinism flags
(\texttt{-fno-fast-math}) enabled by default; the federal-statute build is
pinned to a specific source hash in the \texttt{Cargo.lock}.
\end{remark}

\subsection{Computational Complexity}

Let $T(\mathcal{M})$ denote the wall-clock time for a single METIS call on
$G$ with $k$ districts.
METIS complexity is $O(m \log n)$ in the number of edges and vertices of
$G$~\citep{karypis1998}.

\begin{proposition}[Complexity]
\label{prop:complexity}
In the worst case, \cs{} terminates after $j_{\mathrm{stop}} = j^* + T$ METIS
calls, for a total wall-clock cost of $O((j^* + T) \cdot m \log n)$.
\end{proposition}

\begin{proof}
Each iteration of the while loop calls $\mathcal{M}$ exactly once.
The loop terminates when $\mathtt{no\_improve} = T$, which requires exactly
$j^* + T$ iterations where $j^*$ is the last improving index.
\end{proof}

Empirically (Section~\ref{sec:sufficiency}), $j^* \le 1{,}023$ for all fifty
states in the B.7 dataset, so the practical worst case is
$(1{,}023 + 600) \times T(\mathcal{M}) \approx 1{,}623 \times T(\mathcal{M})$.
For Texas ($n \approx 5{,}000$ tracts, $k = 38$), a single METIS call takes
approximately $0.3$\,s on a modern workstation, making the full sweep
approximately $8.1$\,min---well within the statutory 24-hour window.
Most states are faster by an order of magnitude.

For the practical threshold $\Tprac = 500$, the worst-case empirical cost
is $(1{,}023 + 500) \times T(\mathcal{M})$, and for states where $j^* \le 500$
(all states except Wisconsin in the B.7 dataset), the cost is
$\le 1{,}100 \times T(\mathcal{M})$.
